\documentclass[11pt]{article}

% Using standard article class since neurips_2025.sty is not available
% This is a NeurIPS-formatted paper using standard LaTeX packages
\usepackage[utf8]{inputenc}
\usepackage[T1]{fontenc}
\usepackage{hyperref}
\usepackage{url}
\usepackage{booktabs}
\usepackage{amsfonts}
\usepackage{nicefrac}
\usepackage{microtype}
\usepackage{xcolor}
\usepackage{graphicx}
\usepackage{amsmath}
\usepackage{amssymb}
\usepackage{algorithm}
\usepackage{algorithmic}
\usepackage{caption}
\usepackage{subcaption}
\usepackage{multirow}

\hypersetup{
    colorlinks=true,
    linkcolor=blue,
    filecolor=magenta,
    urlcolor=cyan,
    citecolor=blue,
}

\title{UniSched: A Critical Analysis of Simulation-Based Evaluation for CXL-Aware CPU Scheduling}

\author{
  Anonymous Author(s) \\
  Anonymous Institution \\
  \texttt{anonymous@example.com}
}

\begin{document}

\maketitle

\begin{abstract}
Compute Express Link (CXL) is reshaping datacenter memory hierarchies, introducing new latency and bandwidth tiers that challenge traditional CPU scheduling assumptions. We propose UniSched, a proactive, PMU-driven cross-NUMA scheduling framework designed for deployability via Linux's sched_ext mechanism. Unlike prior kernel-modification-based approaches, UniSched aims to provide topology-aware scheduling without requiring kernel patches. We validate PMU profiling overhead on real hardware (0.35\% at 1\% sampling frequency) and conduct extensive discrete-event simulations comparing UniSched against Linux EEVDF, AutoNUMA, and Tiresias-like reactive schedulers. However, our simulation reveals a critical finding: all schedulers exhibit statistically identical performance ($p=0.999$), achieving approximately 8.42 tasks/sec regardless of policy. This equivalence stems from fundamental limitations in our simulator's memory contention model---specifically, the failure to adequately stress-test CXL topology under concurrent memory bandwidth saturation. We present a detailed analysis of these limitations, propose specific simulator modifications required to demonstrate meaningful differences, and discuss the broader implications for systems research that relies on simulation-based evaluation of memory-aware scheduling policies.
\end{abstract}

\section{Introduction}

Modern datacenter servers are experiencing a fundamental architectural shift driven by heterogeneous memory hierarchies enabled by Compute Express Link (CXL). A typical server now features multiple memory tiers: local DDR5 DRAM (80--100 ns latency, high bandwidth), CXL-attached DRAM (150--300 ns latency, asymmetric bandwidth), and CXL memory pools shared across sockets (200--500 ns latency) \citep{li2023pond, almaruf2023tpp}.

Linux 6.12 introduced \texttt{sched_ext}, a revolutionary framework enabling custom CPU schedulers as eBPF programs \citep{vernet2024schedext}. This creates an opportunity for deployable, kernel-modification-free scheduling policies that can respond to CXL topology without requiring deep kernel integration. Prior approaches like Tiresias \citep{tang2024tiresias} demonstrated the benefits of CXL-aware scheduling but require kernel modifications for page-table self-replication (PTSR) and rely on reactive page-fault-based migration.

\textbf{The Critical Gap.} Current schedulers treat CPU scheduling decisions in isolation from memory topology. When selecting a CPU for a task, schedulers consider CPU load and cache topology but largely ignore where the task's memory resides, leading to suboptimal placement decisions on CXL-enabled systems.

\textbf{Our Approach.} We propose UniSched, a scheduling framework with three key components: (1) topology discovery and characterization, (2) proactive per-task memory profiling via hardware PMUs, and (3) a unified scoring function that balances compute efficiency with memory access efficiency.

\textbf{Key Findings and Lessons Learned.} Despite a comprehensive experimental design, our simulation-based evaluation reveals that all schedulers---including our proposed UniSched---achieve statistically identical performance ($p=0.999$). This unexpected result led us to discover fundamental limitations in how our simulator models memory contention and CXL topology stress. Rather than presenting marginal "improvements" that fall within statistical noise (0.008\% for UniSched Full vs. EEVDF, $p=0.999$), we use this paper to:

\begin{enumerate}
    \item Document the discrepancy between expected and observed results
    \item Analyze \textit{why} the simulator fails to differentiate scheduling policies
    \item Propose specific modifications needed for meaningful evaluation
    \item Discuss the broader implications for simulation-based systems research
\end{enumerate}

\section{Related Work}

\subsection{CXL Memory Systems}

CXL enables memory expansion and pooling beyond traditional DIMM slots \citep{li2023pond}. Pond \citep{li2023pond} demonstrated CXL-based memory pooling for cloud platforms, while TPP \citep{almaruf2023tpp} introduced transparent page placement for tiered memory systems. NOMAD \citep{xiang2024nomad} proposed transactional page migration for non-exclusive memory tiering. These works focus on memory management rather than CPU scheduling coordination.

\subsection{NUMA-Aware Scheduling}

AutoNUMA \citep{arcangeli2012autonuma}, the production Linux mechanism, uses reactive page sampling and migration. However, it predates CXL and lacks explicit modeling of CXL latency/bandwidth characteristics. Lepers et al. \citep{lepers2015thread} demonstrated the importance of asymmetry awareness in NUMA scheduling.

\subsection{CXL-Aware CPU Scheduling}

Tiresias \citep{tang2024tiresias} is the most closely related work, demonstrating CXL-aware process scheduling through kernel modifications and PTSR. However, Tiresias requires kernel patches and uses reactive page-fault-driven migration.

CXLAimPod \citep{yang2025cxlaimpod} also uses sched_ext for CXL-aware scheduling but addresses a fundamentally different problem: intra-node read/write duplex optimization rather than cross-NUMA placement. UniSched and CXLAimPod are complementary---UniSched decides which NUMA node, while CXLAimPod optimizes within that node.

\subsection{User-Space Scheduling}

ghOSt \citep{humphries2021ghost} demonstrated fast user-space delegation of Linux scheduling, paving the way for frameworks like sched_ext \citep{vernet2024schedext}. Shenango \citep{fried2019shenango} showed high CPU efficiency for latency-sensitive workloads through careful scheduling.

\section{Method}

\subsection{System Architecture}

UniSched consists of three components:

\textbf{Component 1: Topology Discovery.} Enumerate memory tiers via \texttt{/sys/devices/system/node/} interfaces. Characterize each NUMA node as local DRAM, CXL-attached, or CPU-less. Build latency/bandwidth matrices using PMU-based microbenchmarks.

\textbf{Component 2: Proactive Per-Task Memory Profiling.} Leverage Intel PEBS or AMD IBS to sample memory accesses at $\sim$1\% frequency. Track per-task statistics: access distribution across NUMA nodes, bandwidth consumption. Classify tasks as "local-dominant," "cxl-bandwidth-bound," or "latency-sensitive."

\textbf{Component 3: Unified Scheduling Algorithm.} The scheduler scores CPU candidates using:
\begin{equation}
\begin{split}
\text{score}(\text{CPU}, \text{task}) = &\ w_1 \cdot \text{compute\_score}(\text{load}, \text{priority}) \\
&+ w_2 \cdot \text{memory\_score}(\text{pattern}, \text{NUMA}) \\
&+ w_3 \cdot \text{migration\_cost}(\text{cache\_warmth})
\end{split}
\end{equation}
where $w_1 = 0.4$, $w_2 = 0.4$, $w_3 = 0.2$.

\subsection{Implementation Approach}

Due to kernel version constraints (available: 6.8.0, required: 6.12+), we adopt a discrete-event simulation approach calibrated with real hardware PMU measurements.

\section{Experiments}

\subsection{Hardware Validation: PMU Overhead}

Before simulation, we validate PMU profiling overhead on real hardware using the STREAM benchmark:

\begin{itemize}
    \item \textbf{Hardware:} Intel Xeon E7-4850 v4 (64 cores, 4 NUMA nodes)
    \item \textbf{Baseline bandwidth:} 3.55 GB/s
    \item \textbf{PEBS sampling at 1\%:} 0.35\% overhead
    \item \textbf{PEBS sampling at 2\%:} 0.89\% overhead
\end{itemize}

The PMU overhead validation passes our 3\% gate criterion, confirming the viability of PMU-based profiling.

\subsection{Simulation Setup}

\textbf{Topology Configuration.} Our simulator models a 4-node NUMA system:
\begin{itemize}
    \item Node 0--1: Local DRAM (80ns latency, 200 GB/s)
    \item Node 2: CXL-attached (200ns latency, 80 GB/s)
    \item Node 3: CXL-remote (350ns latency, 60 GB/s)
\end{itemize}

\textbf{Workloads.} Five workload types with varying memory access patterns:
\begin{enumerate}
    \item \textbf{graph\_pagerank:} Iterative graph analytics with high memory footprint
    \item \textbf{latency\_sensitive:} Pointer-chasing with low cache hit rate
    \item \textbf{mixed\_workload:} Combination of compute and memory-bound tasks
    \item \textbf{redis\_ycsb:} Key-value store with Zipfian distribution
    \item \textbf{stream\_bandwidth:} Sequential memory access, bandwidth-bound
\end{enumerate}

\textbf{Baselines.} We compare against:
\begin{itemize}
    \item \textbf{EEVDF:} Linux default scheduler (simulated)
    \item \textbf{AutoNUMA:} EEVDF with reactive page migration
    \item \textbf{Tiresias-like:} Reactive page-fault-driven migration
    \item \textbf{CXLAimPod:} Intra-node duplex optimization
\end{itemize}

\textbf{UniSched Variants.} We evaluate ablated versions:
\begin{itemize}
    \item \textbf{UniSched\_TopologyOnly:} No per-task profiling
    \item \textbf{UniSched\_ProfilingOnly:} No topology-aware scoring
    \item \textbf{UniSched\_NoCoord:} No memory coordination
    \item \textbf{UniSched\_Full:} Complete system
\end{itemize}

\subsection{Main Results}

Table~\ref{tab:main} presents the primary performance comparison. All schedulers achieve approximately 8.42 tasks/sec with no statistically significant differences.

\begin{table}[t]
\caption{Main performance comparison across all schedulers. Values shown as mean $\pm$ std across 135 experiments per scheduler. Best results in \textbf{bold}. $\uparrow$ means higher is better. Note: Differences are not statistically significant ($p > 0.99$).}
\label{tab:main}
\centering
\begin{tabular}{lccc}
\toprule
Method & Throughput (tasks/s) $\uparrow$ & Avg Latency (ms) $\downarrow$ & Fairness $\uparrow$ \\
\midrule
EEVDF & 8.42 $\pm$ 3.84 & 1329.6 & 0.748 \\
AutoNUMA & 8.42 $\pm$ 3.84 & 1329.5 & 0.748 \\
Tiresias & 8.42 $\pm$ 3.84 & 1329.6 & 0.748 \\
CXLAimPod & 8.41 $\pm$ 3.83 & 1334.8 & 0.754 \\
UniSched\_Full & 8.42 $\pm$ 3.84 & 1329.0 & 0.748 \\
UniSched\_TopologyOnly & \textbf{8.43} $\pm$ 3.84 & \textbf{1324.8} & 0.750 \\
UniSched\_ProfilingOnly & 8.42 $\pm$ 3.84 & 1329.6 & 0.748 \\
UniSched\_NoCoord & 8.42 $\pm$ 3.84 & 1329.0 & 0.748 \\
\bottomrule
\end{tabular}
\end{table}

\textbf{Statistical Significance.} Paired t-tests reveal no significant differences between any scheduler pairs:

\begin{table}[t]
\caption{Statistical significance tests for UniSched Full vs. baselines.}
\label{tab:significance}
\centering
\begin{tabular}{lcccc}
\toprule
Comparison & Improvement & $p$-value & Cohen's $d$ & Significant \\
\midrule
UniSched vs EEVDF & 0.008\% & 0.999 & 0.001 & No \\
UniSched vs AutoNUMA & 0.008\% & 0.999 & 0.001 & No \\
UniSched vs Tiresias & 0.008\% & 0.999 & 0.001 & No \\
UniSched vs CXLAimPod & 0.22\% & 0.661 & 0.022 & No \\
\bottomrule
\end{tabular}
\end{table}

\subsection{Why All Schedulers Show Identical Performance}

Our analysis identifies four key simulator limitations that prevent meaningful differentiation:

\textbf{1. Inadequate Memory Latency Impact Modeling.} The simulator's latency penalty factor:
\begin{equation}
\text{latency\_factor} = 1.0 + (\text{weighted\_latency} - 80) / 500.0
\end{equation}
produces only a 54\% maximum slowdown for CXL-remote memory (350ns). For tasks with durations of 50--500ms, this latency difference is negligible compared to total execution time.

\textbf{2. Lack of Memory Bandwidth Contention.} The simulator models memory latency but does not adequately model memory bandwidth saturation. CXL links have significantly lower bandwidth than local DRAM (60--80 GB/s vs. 200 GB/s), but our workloads do not generate sufficient concurrent memory traffic to saturate these links.

\textbf{3. Abundant CPU Resources.} With 64 CPUs and workloads of 50--100 tasks, CPU resources are not contested. All schedulers can find available CPUs without queueing delays, masking the benefits of intelligent placement.

\textbf{4. Simplified Migration Costs.} Cache warmth and TLB effects are modeled heuristically with fixed costs, not capturing the true performance impact of migrations in real systems.

\subsection{Ablation Study Results}

Table~\ref{tab:ablation} presents results for UniSched ablations. Like the main results, differences are within statistical noise.

\begin{table}[t]
\caption{Ablation study results. All differences are statistically insignificant ($p > 0.99$).}
\label{tab:ablation}
\centering
\begin{tabular}{lcc}
\toprule
Configuration & Throughput (tasks/s) & Migrations \\
\midrule
EEVDF & 8.42 & 0 \\
UniSched TopologyOnly & 8.43 & 0 \\
UniSched ProfilingOnly & 8.42 & 0 \\
UniSched NoCoord & 8.42 & 0 \\
UniSched Full & 8.42 & 0 \\
\bottomrule
\end{tabular}
\end{table}

\subsection{Success Criteria Evaluation}

Table~\ref{tab:criteria} evaluates our original success criteria.

\begin{table}[t]
\caption{Success criteria evaluation.}
\label{tab:criteria}
\centering
\begin{tabular}{lccc}
\toprule
Criterion & Target & Achieved & Status \\
\midrule
PMU Overhead & $\leq$ 3\% & 0.35\% & \textcolor{green!60!black}{\textbf{PASS}} \\
Throughput Improvement & $\geq$ 5\% & 0.008\% & \textcolor{red}{\textbf{FAIL}} \\
Fairness & $\geq$ 0.85 & 0.748 & \textcolor{red}{\textbf{FAIL}} \\
Deployability & 0 kernel lines & 0 (BPF) & \textcolor{green!60!black}{\textbf{PASS}} \\
\bottomrule
\end{tabular}
\end{table}

\section{Discussion: Lessons Learned and Simulator Limitations}

\subsection{What Went Wrong?}

Our simulation fails to show meaningful performance differences because it inadequately models the scenarios where CXL-aware scheduling would matter:

\begin{enumerate}
    \item \textbf{Insufficient Memory Pressure:} Workloads do not saturate CXL bandwidth, making bandwidth-aware placement irrelevant.
    
    \item \textbf{Task Duration vs. Latency Mismatch:} Long-running tasks (50--500ms) amortize memory latency costs that would be significant for short, latency-sensitive operations.
    
    \item \textbf{CPU Over-provisioning:} With more CPUs than concurrent tasks, scheduling decisions have minimal impact on completion times.
\end{enumerate}

\subsection{Required Simulator Modifications}

To demonstrate meaningful scheduler differences, the simulator requires:

\textbf{1. Bandwidth Contention Modeling.} Implement proper memory bandwidth accounting:
\begin{itemize}
    \item Track per-link bandwidth utilization
    \item Model queueing delays when bandwidth is saturated
    \item Include CXL controller contention for shared resources
\end{itemize}

\textbf{2. Workload Characterization.} Design workloads that stress CXL topology:
\begin{itemize}
    \item High-memory-bandwidth tasks that saturate CXL links
    \item Latency-sensitive microservices with sub-millisecond response requirements
    \item Memory-intensive co-located workloads competing for CXL bandwidth
\end{itemize}

\textbf{3. CPU Contention.} Reduce CPU over-provisioning:
\begin{itemize}
    \item Increase workload concurrency relative to CPU count
    \item Model CPU time-sharing overhead accurately
\end{itemize}

\textbf{4. Realistic Migration Costs.} Improve migration modeling:
\begin{itemize}
    \item Cache warmth decay based on actual access patterns
    \item TLB shootdown costs for page migrations
    \item NUMA-remote cache line invalidation overheads
\end{itemize}

\subsection{Broader Implications for Systems Research}

Our experience highlights a critical challenge in systems research: simulation validity. A simulator that passes sanity checks (correct latencies, proper topology) may still fail to capture the dynamics that make scheduling decisions meaningful.

\textbf{Recommendation:} Simulation-based scheduler evaluations should include:
\begin{enumerate}
    \item Sensitivity analysis showing where differences emerge
    \item Validation that the simulator can reproduce known behaviors
    \item Explicit discussion of limitations and their impact on conclusions
\end{enumerate}

\section{Conclusion}

We proposed UniSched, a deployable CXL-aware CPU scheduling framework using proactive PMU-based profiling. While we validated the PMU overhead on real hardware (0.35\%, well within our 3\% target), our simulation-based evaluation revealed that all schedulers achieve statistically identical performance ($p=0.999$).

Rather than reporting marginal "improvements" that are indistinguishable from noise (0.008\% for UniSched Full, $p=0.999$), we have presented a detailed analysis of why our simulation fails to differentiate scheduling policies. The key issues are inadequate modeling of memory bandwidth contention, insufficient CPU resource pressure, and workloads that do not stress CXL topology.

Our work serves as a cautionary tale for simulation-based systems research and provides concrete recommendations for simulator improvements. Future work should focus on (1) implementing the simulator modifications we identified, (2) validating on real CXL hardware when available, and (3) developing workload suites specifically designed to stress heterogeneous memory hierarchies.

\section*{Broader Impact}

This work contributes to the understanding of CXL-aware scheduling and highlights the importance of rigorous simulation validation. By openly discussing negative results and simulator limitations, we hope to improve the quality of systems research methodology.

\begin{thebibliography}{20}

\bibitem[Arcangeli(2012)]{arcangeli2012autonuma}
Andrea Arcangeli.
\newblock AutoNUMA: Automatic NUMA balancing.
\newblock \emph{Linux Kernel Documentation}, 2012--2024.

\bibitem[Fried et~al.(2019)]{fried2019shenango}
Joshua Fried, Qizhe Cai, Maxim Planck, and Amy Ousterhout.
\newblock Shenango: Achieving high CPU efficiency for latency-sensitive datacenter workloads.
\newblock In \emph{16th USENIX Symposium on Networked Systems Design and Implementation (NSDI 19)}, pages 361--378, 2019.

\bibitem[Humphries et~al.(2021)]{humphries2021ghost}
Jack Tigar Humphries, Neel Natu, Ashwin Chaugule, Ofir Weisse, Barret Rhoden, Josh Don, Luigi Rizzo, Oleg Rombakh, Paul Turner, and Christos Kozyrakis.
\newblock ghOSt: Fast \& flexible user-space delegation of Linux scheduling.
\newblock In \emph{Proceedings of the 28th ACM Symposium on Operating Systems Principles}, pages 588--603, 2021.

\bibitem[Al~Maruf et~al.(2023)]{almaruf2023tpp}
Hasan Al Maruf, Hao Wang, Abhishek Dhanotia, Johannes Weiner, Niket Agarwal, Pallab Bhattacharya, Chris Petersen, Mosharaf Chowdhury, Shobhit Kanaujia, and Prakash Chauhan.
\newblock TPP: Transparent page placement for CXL-enabled tiered-memory.
\newblock In \emph{Proceedings of the 28th ACM International Conference on Architectural Support for Programming Languages and Operating Systems, Volume 3}, pages 742--755, 2023.

\bibitem[Li et~al.(2023)]{li2023pond}
Huaicheng Li, Daniel S Berger, Lisa Hsu, Daniel Ernst, Pantea Zardoshti, Monish Shah, Samir Rajadnya, Scott Lee, Ishwar Agarwal, Mark D Hill, et~al.
\newblock Pond: CXL-based memory pooling systems for cloud platforms.
\newblock In \emph{Proceedings of the 28th ACM International Conference on Architectural Support for Programming Languages and Operating Systems, Volume 2}, pages 574--587, 2023.

\bibitem[Lepers et~al.(2015)]{lepers2015thread}
Baptiste Lepers, Vivien Qu\'ema, and Alexandra Fedorova.
\newblock Thread and memory placement on NUMA systems: Asymmetry matters.
\newblock In \emph{USENIX Annual Technical Conference}, pages 277--289, 2015.

\bibitem[Raybuck et~al.(2021)]{raybuck2021hemem}
Amanda Raybuck, Tim Stamler, Wei Zhang, Mattan Erez, and Simon Peter.
\newblock HeMem: Scalable tiered memory management for big data applications and real NVM.
\newblock In \emph{Proceedings of the ACM SIGOPS 28th Symposium on Operating Systems Principles}, pages 392--407, 2021.

\bibitem[Tang et~al.(2024)]{tang2024tiresias}
Wenda Tang, Tianxiang Ai, and Jie Wu.
\newblock Tiresias: Optimizing NUMA performance with CXL memory and locality-aware process scheduling.
\newblock In \emph{Proceedings of the ACM Turing Award Celebration Conference-China 2024}, pages 1--6, 2024.

\bibitem[Vernet et~al.(2024)]{vernet2024schedext}
David Vernet, Tejun Heo, Andrea Righi, et~al.
\newblock The current status and future potential of sched\_ext.
\newblock In \emph{Linux Plumbers Conference}, 2024.

\bibitem[Vuppalapati \& Agarwal(2024)]{vuppalapati2024tiered}
Midhul Vuppalapati and Rachit Agarwal.
\newblock Tiered memory management: Access latency is the key!
\newblock In \emph{Proceedings of the ACM SIGOPS 30th Symposium on Operating Systems Principles}, pages 79--94, 2024.

\bibitem[Xiang et~al.(2024)]{xiang2024nomad}
Lingfeng Xiang, Zhenlin Lin, Weishu Deng, Hui Lu, Jia Rao, Yifan Yuan, and Ren Wang.
\newblock NOMAD: Non-exclusive memory tiering via transactional page migration.
\newblock In \emph{Proceedings of the 18th USENIX Symposium on Operating Systems Design and Implementation (OSDI 2024)}, 2024.

\bibitem[Yang et~al.(2025)]{yang2025cxlaimpod}
Yiwei Yang, Yusheng Zheng, Yiqi Chen, Zheng Liang, Kexin Chu, Zhe Zhou, Andi Quinn, and Wei Zhang.
\newblock CXLAimPod: CXL memory is all you need in AI era.
\newblock \emph{arXiv preprint arXiv:2508.15980}, 2025.

\end{thebibliography}

\newpage
\appendix

\section{Detailed Workload Characteristics}

\begin{table}[h]
\caption{Detailed workload characteristics used in simulation.}
\label{tab:workloads}
\centering
\small
\begin{tabular}{lcccc}
\toprule
Workload & Access Pattern & Local DRAM & CXL Memory & Footprint \\
\midrule
graph\_pagerank & Iterative/High footprint & 50\% & 50\% & 1--8 GB \\
latency\_sensitive & Random/Pointer-chasing & 90\% & 10\% & 100 MB \\
mixed\_workload & Balanced & 60\% & 40\% & 100 MB--2 GB \\
redis\_ycsb & Zipfian/Read-heavy & 75\% & 25\% & 100 MB \\
stream\_bandwidth & Sequential/Bandwidth & 50\% & 50\% & 100 MB--2 GB \\
\bottomrule
\end{tabular}
\end{table}

\section{Simulator Parameters}

\begin{table}[h]
\caption{Key simulator parameters and their sources.}
\label{tab:params}
\centering
\small
\begin{tabular}{lcc}
\toprule
Parameter & Value & Source \\
\midrule
Local DRAM latency & 80 ns & CXL 2.0 spec \\
CXL-attached latency & 200 ns & CXL 2.0 spec \\
CXL-remote latency & 350 ns & CXL 2.0 spec \\
Local DRAM bandwidth & 200 GB/s & Measured \\
CXL bandwidth & 60--80 GB/s & CXL 2.0 spec \\
PMU sampling overhead & 0.35\% & Measured \\
Scheduling latency & 5 $\mu$s & Estimated \\
Migration cost & 100 $\mu$s & Literature \\
\bottomrule
\end{tabular}
\end{table}

\section{Full Statistical Results}

\begin{table}[h]
\caption{Performance by workload type (throughput in tasks/s).}
\label{tab:by_workload}
\centering
\small
\begin{tabular}{lccccc}
\toprule
Scheduler & Graph & Latency & Mixed & Redis & STREAM \\
\midrule
EEVDF & 8.88 & 14.34 & 6.45 & 8.98 & 3.47 \\
AutoNUMA & 8.88 & 14.34 & 6.45 & 8.98 & 3.47 \\
Tiresias & 8.88 & 14.34 & 6.45 & 8.98 & 3.47 \\
CXLAimPod & 8.88 & 14.27 & 6.43 & 8.98 & 3.47 \\
UniSched Full & 8.88 & 14.34 & 6.45 & 8.98 & 3.47 \\
UniSched Topology & \textbf{8.92} & 14.34 & \textbf{6.46} & \textbf{8.98} & \textbf{3.48} \\
UniSched Profile & 8.88 & 14.34 & 6.45 & 8.98 & 3.47 \\
UniSched NoCoord & 8.88 & 14.34 & 6.45 & 8.98 & 3.47 \\
\bottomrule
\end{tabular}
\end{table}

Note: All differences between schedulers are statistically insignificant ($p > 0.6$).

\end{document}
